\scanpage{37}
\mspage{36}

Il contient de plus les \deleted{objets} figures élémentaires, i.e. satisfait :
\begingroup\makeatletter\tagsleft@true\makeatother
\[
 \forall\,\unsure{x|X}\in M,\quad \{\{X\}\}\in\unsure{\mathcal F|F}
 \tag*{(iii)}
\]
\endgroup
et par suite, il contient toutes les parties fermées de réunions finies de tels $\widetilde X$.

Inversement, \unsure{tte} $\mathcal F$ satisfaisant ces trois conditions satisfait aux conditions voulues

\deleted{Si \unsure{les}} \inserted{Pour que les} figures de $\mathcal F$ soient de \unsure{t.f.}, \unsure{il f} et s. qu'on \deleted{ait}
\begingroup\makeatletter\tagsleft@true\makeatother
\[
 \mathcal F=\mathcal F_{0}=
 \left\{\begin{aligned}
 &\text{réunions finies d'ensembles}\\
 &\qquad\text{de la forme }\widetilde X
 \end{aligned}\right\}
 \tag*{a)}
\]
\endgroup

\unsure{(ce|Ce} qui ne donne un atlas que si on suppose de plus

b) toute partie fermée d'un $\widetilde X$ \inserted{(pour $\leq$)} est de type fini

À part donc un petit flottement dû à la possibilité de \deleted{\unsure{feuilles|feuillets}} \inserted{figures} \deleted{\unsure{finies}} infinies, \deleted{\unsure{o}}\unsure{on} voit que \unsure{tout|\deleted{tout}} s'exprime par $(M,{\leq},{\ll},{\parallel})$, avec $\parallel$ \unsure{définie|défini} sur $\unsure{L|\mathcal L}$ sans plus, \deleted{\unsure{(et si on dit aussi|(et sinon de aussi|(et sinon de deux}} satisfaisant aux quatre axiomes (cinq avec $(\mathrm{ML\,0})$) $\mathrm{ML\,1}-\mathrm{ML\,4}$ (au lieu
